\cotefile{157-3}
\scanpage{114}
% line head | text | 63
\mspage{63}
% line 1 | text | donne ceci
\unsure{donne} ceci :

% line 2–10 | text | Proposition 2
\uline{Proposition 2} : Considérons une situation de recollement
de deux topos $\mathscr U$, $\mathscr T$ par $\varphi:\operatorname{Fais}\mathscr U\unsure{-|\to}\operatorname{Fais}\mathscr T$,
foncteur exact à gauche. Supposons que $\mathscr U=\operatorname{Top}(U)$,
$\mathscr T=\operatorname{Top}(T)$, avec $U$ et $T$ dans $\mathrm{Cat}$. Pour que
le topos recollé $\unsure{\mathscr X|\mathscr Z}$ soit équivalent à un topos
$\operatorname{Top}(X)$, avec $X$ dans $\mathrm{Cat}$, il f. et s. que
le foncteur de recollement $\varphi$ commute aux
$\lim_{\leftarrow}$ \unsure{(quelc.),|(quelc.,)} ou ce qui revient au $=$,
qu'il ait un adjoint à gauche :
% line 11 | display | ψ, Fais T, Fais U, the two presheaf categories
\[
\psi:\underset{\substack{\shortparallel\\T^\wedge}}{\underline{\operatorname{Fais}\mathscr T}}
\longrightarrow\underset{\substack{\shortparallel\\U^\wedge}}{\underline{\operatorname{Fais}\mathscr U}} .
\]

% line 12–16 | text | On s'attend dès lors
On s'attend dès lors que \deleted{\unsure{cette}} la donnée de
recollement $\varphi$, ou ce qui revient au $=$ le
foncteur $\psi$ \unsure{a}\deleted{\unsure{yant}} un adjoint à droite, i.e.
commutant aux $\lim_{\rightarrow}$ quelc., revienne à
celle d'une donnée de recollement $\unsure{H|M}$
% line 17 | display | H, U with a small circle, T, Ens
\[
H:U^{\mathrm o}\times T\longrightarrow\mathrm{Ens}
\]
% line 18–24 | text | pour U et T
pour $U$ et $T$. Mais on sait bien que
les foncteurs $\psi:T^\wedge\longrightarrow\mathscr C$ dans une
catégorie \inserted{$\mathscr C$} avec $\lim_{\rightarrow}$ quelc., qui commutent
aux dites $\lim_{\rightarrow}$, forment une catégorie
équivalente (\unsure{par|pour} le foncteur $\psi\longmapsto\psi\mid T\,(\subset T^\wedge)$)
à celle des \deleted{\unsure{préf.}} foncteurs
% line 25 | display | ψ₀, T, C, U with wedge, Hom
\[
\psi_0:T\longrightarrow\mathscr C
\quad \bigl(=\underset{\substack{\shortparallel\\\underline{\operatorname{Hom}}(\deleted{\unsure{U}}U^{\mathrm o},\mathrm{Ens})}}{U^\wedge}
\text{ en l'occurrence}\bigr).
\]
% archive sheet number: 114
\scanpage{115}
% line head | text | 64
\mspage{64}
% line 1–2 | text | et la donnée d'un tel foncteur
et \inserted{la donnée d'}un tel foncteur $\psi_0$ revient bien à celle d'un
$H$. Donc

% line 3–13 | text | Corollaire. La catégorie des données
\uline{Corollaire}\quad La catégorie des données de recollement
topossiques $\varphi:\deleted{\unsure{T}}U^\wedge\to T^\wedge$ entre $\operatorname{Top}(\deleted{\unsure{T}}U),\operatorname{Top}(T)$,
telles que \deleted{$\varphi$ soit} le topos recollé soit \inserted{un ``topos \uline{\unsure{paradisiaque}}'', i.e. \unsure{soit}} équivalent à
un topos de la forme $X^\wedge$ (avec $X$ dans $\mathrm{Cat}$),
est équivalente à celle des données de recollement
$H:U^{\mathrm o}\times T\unsure{-|\to}\mathrm{Ens}$ entre $U$ et $T$. \deleted{\unsure{Elle}} La
\deleted{\unsure{pour un}} donnée \inserted{topossique} $\varphi$ associée à $H$ est le foncteur adjoint
à droite du $T^\wedge\longrightarrow U^\wedge$ déduit de $T\xrightarrow{H}U^\wedge$
par prolongement commutant aux $\lim_{\rightarrow}$, ou encore
elle est donnée par
% line 14 | display | φ(F)(t), Hom on U with wedge
\[
\varphi(\underset{\substack{\rotatebox{-90}{$\in$}\\U^\wedge}}{F})(t)
\simeq\operatorname{Hom}_{U^\wedge}(\deleted{\unsure{F},}H(t),F),
\]
% line 15–19 | text | où dans le second membre
où dans le second membre $H$ est interprété comme
un foncteur $T\longrightarrow U^\wedge$. Le topos recollé
\inserted{par $\varphi$} est can. équivalent à $\operatorname{Top}(X)$, où $X$ est
l'objet de $\mathrm{Cat}$ obtenu par recollement à l'aide
de $H$.

% line 20–22 | text | Du point de vue de la théorie des faisceaux
Du point de vue de la théorie des faisceaux,
cela donne l'interprétation de $H$, comme
exprimant le foncteur
% line 23 | display | φ equals j upper star i lower star
\[
\varphi=j^*i_*:U^\wedge-T^\wedge
\]
% line 24 | text | par la formule précédente
par la formule précédente.
% archive sheet number: 115
\scanpage{116}
% line head | text | 65
\mspage{65}
% line 1–2 | text | 1.3. Flèches dans Cat
\textbf{1.3.}\quad \deleted{\unsure{Foncteurs}} Flèches dans $\mathrm{Cat}$, de source ou but un objet
recollé $(X,U,T)$.

% line 3–4 | text | Quand (X,U,T) est donnée
Quand $(X,U,T)$ est donnée par \unsure{leur|leurs} donnée de
recollement
% line 5 | display | U, T, H
\[
U,T,\qquad H:U^{\mathrm o}\times T\longrightarrow\mathrm{Ens},
\]
% line 6 | text | les flèches
les \deleted{foncteurs} flèches
% line 7 | display | f, X prime, X
\[
f:X'\longrightarrow X
\]
% line 8–9 | text | dans Cat se décrivent
dans $\mathrm{Cat}$ se décrivent immédiatement \unsure{ainsi} : \deleted{\unsure{pour}}
la donnée de $f$ définit d'abord
% line 10 | display | U prime and T prime
\[
U'=\unsure{j|f}^{-1}(U),\qquad T'=f^{-1}(T)
\]
% line 11–13 | text | couple d'un ouvert
couple d'un ouvert de $X'$ et du fermé complémentaire,
\deleted{\unsure{et}} de sorte que cette \deleted{\unsure{d.}} donnée équivaut
à celle de

% line 14 | text | 1°) un couple
$\underline{1}^{\circ})$ un couple $(U',T')$ dans $X'$

% line 15–17 | text | 2°) un morphisme
$2^{\circ})$ un morphisme (de $\unsure{\mathscr T|\mathcal T}$) de $(X',U',T')$
\unsure{dans} $(X,U,T)$, ce qui équivaut par 1.1.
à un \unsure{triple}
% line 18–20 | display | f on U prime, f on T prime, u
\[
\left\{\begin{aligned}
&f_{U'}:U'\to U\\
&f_{T'}:T'\unsure{-|\to}T\\
&\unsure{u|v}:H_{U',T'}\longrightarrow H\circ(f_{U'}^{\mathrm o}\times f_{T'})
\end{aligned}\right.
\]
% line 21 | text | où on a posé
où on a posé
% line 22–23 | display | H with subscripts, its value at u prime and t prime
\[
\left\{\begin{aligned}
&H_{U',T'}:{U'}^{\mathrm o}\times T'\longrightarrow\mathrm{Ens}\\
&H_{U',T'}(u',t')=\operatorname{Hom}_{X'}(u',t'),
\end{aligned}\right.
\]
% line 24–25 | text | i.e. H is the gluing datum
i.e. $H_{U',T'}$ est la donnée de recollement
qui correspond au triple $(X',U^{\unsure{\prime\prime|\prime}},T')$ de $\mathscr T$.
% archive sheet number: 116
\scanpage{117}
% line head | text | 66
\mspage{66}
% line 1–3 | text | Nous serons plus intéressés
\unsure{Nous} serons plus intéressés par la \deleted{\unsure{problème}}\inserted{situation} inverse,
\deleted{\unsure{d'écrire}} par la question de déterminer les
foncteurs
% line 4 | display | X, X subscript U,T,H, X prime
\[
\underset{\substack{\shortparallel\\X_{U,T,H}}}{X}\longrightarrow X'
\]
% line 5–7 | text | de X, obtenu par
de $X$, obtenu par $(U,T,H)$, dans un objet
quelconque de $\mathrm{Cat}$. On trouve tautologiquement

% line 8–10 | text | Proposition 3
\uline{Proposition 3}\quad $\operatorname{Hom}_{\mathrm{Cat}}(X,X')$ est en corr. bijective
canonique, fonctorielle en $X'$ dans $\mathrm{Cat}$, avec
l'ens des triples
% line 11–14 | display | f_U, f_T, u and the functor category
\[
\left\{\begin{aligned}
&\left.\begin{aligned}
f_U &: U\to X'\\
f_T &: T\to X'
\end{aligned}\right\}\ \text{flèches \unsure{dans|de} $\mathrm{Cat}$}\\
&u:H\longrightarrow H_{X'}(f_U^{\mathrm o}\times f_T)\quad\text{flèche \unsure{dans|de}}\\
&\hspace{9em}\underline{\operatorname{Hom}}(U^{\mathrm o}\times T,\mathrm{Ens})
\end{aligned}\right.
\]
% line 15 | text | où pour tt X prime
où pour \unsure{tt} $X'$ dans $\mathrm{Cat}$, on a posé
% line 16 | display | H of X prime and Hom
\[
H_{X'}=\bigl((x,y)\unsure{\mapsto|\vdash}\operatorname{Hom}_{X'}(x,y)\bigr):
{X'}^{\mathrm o}\times X'\longrightarrow\mathrm{Ens}.
\]

% line 17–18 | text | Sûrement on aura
\unsure{Sûrement} on \unsure{aura} $=$ un \deleted{équivalence}\inserted{isomorphisme} de
catégories
% line 19–20 | display | Underlined Hom and category of triples
\[
\begin{aligned}
\underline{\operatorname{Hom}}(X,X')&\simeq\text{catégorie des triples }(f_U,f_T,u)\\[-2pt]
&\hspace{3em}\text{comme \unsure{dans $(*)$|dessus}}.
\end{aligned}
\]

% line 21–22 | text | 1.4. Construction de sommes amalgamées
\uline{1.4}.\quad \uline{Construction de sommes amalgamées} \unsure{à} partir
\uline{d'une immersion ouverte ou fermée}.

% line 23–24 | text | Pour une construction
Pour une construction d'une \inserted{sorte de} somme
amalgamée, en termes d'un triple $(X,U,T)+\ldots$,
% archive sheet number: 117
\scanpage{118}
% line head | text | 67
\mspage{67}
% line 1–3 | text | qui soit symétrique en U, T
qui soit symétrique en $U,T$, j'ai envie
de regarder d'abord la situation ``mixte'', d'un
diagramme en traits pleins
% line 4–6 | diagram | (α) U to X from T over U prime to X prime from T prime
\diagram{\begin{tikzcd}[ampersand replacement=\&,column sep=small]
\& U\arrow[r,hook,"i"]\arrow[d,"f_U"'] \& X\arrow[d,dashed,"f"] \& T\arrow[l,hook',"j"']\arrow[d,"f_T"]\\
{(\unsure{\alpha|a})} \& U'\arrow[r,dashed,hook,"i'"] \& X' \& T'\arrow[l,dashed,hook',"j'"'] {,}
\end{tikzcd}}
% line 7–14 | text | qu'on aimerait compléter
qu'on aimerait compléter ``canoniquement'', si faire
se peut, en un morphisme de triples (i.e. en
une flèche de $\unsure{\mathscr T|\mathcal T}$). L'idée naturelle,
c'est que le $X'$ qui s'y \unsure{insère} \unsure{devrait|serait}
la \unsure{limite} \unsure{inductive}, dans $\mathrm{Cat}$, du diagramme
initial. Si \inserted{$T'=T$,} $f_T=\mathrm{id}_T$, cela signifie donc
que $X'$ s'identifie à la somme amalgamée
$X\amalg_{U\unsure{\alpha|{,}\alpha}}U'$, et inversement si $U'=U$, $f_U=\mathrm{id}_U$.

% line 15–18 | text | Désignons par X prime
Désignons par $X'$ la somme amalgamée et
\deleted{\unsure{défin.}} explicitons, pour $X''$ quelconque \unsure{de|dans} $\mathrm{Cat}$,
\inserted{l'ens} des flèches \unsure{?} de $X'$ dans $X''$, \deleted{\unsure{celui-ci}} comme
foncteur en $X''$. C'est l'ens des triples
% line 19–21 left | diagram | U and T over U prime and T prime, X mapping to X double prime
\diagram{\begin{tikzcd}[ampersand replacement=\&,column sep=small]
 U\arrow[r,hook,"i"]\arrow[d,"f_U"'] \& X\arrow[dd,"g_X"'] \& \deleted{\unsure{U}}T\arrow[l,hook',"j"']\arrow[d,"f_T"]\\
 U'\arrow[dr,"g_{U'}"'] \& \& T'\arrow[dl,"g_{T'}"]\\
 \& X'' \&
\end{tikzcd}}
% line 19–21 right | display | (*) triple g_X, g_U prime, g_T prime and equations
\[
\begin{aligned}
&(g_X,g_{U'},g_{T'})\quad
\left|\begin{aligned}
&g_X:X\to X''\\
&g_{U'}:U'\to X''\\
&g_{T'}:T'\to X''
\end{aligned}\right.\qquad\text{avec}\quad
\begin{aligned}
g_X\circ i&=g_{U'}f_U\\
g_X\circ j&=g_{T'}f_T
\end{aligned}
\end{aligned}\tag*{(*)}
\]
% line 22–25 | text | Notons que pour un tel triple
Notons que pour un tel triple,
$g_X$ est partiellement connu \inserted{déjà} en termes
de $g_{U'},g_{T'}$, par ses restrictions \inserted{$g_X\circ i$ et $g_X\circ j$} à
$U$ et $T$ respectivement, par les formules
% line 26 | display | g_U, g_T
\[
g_U=g_{U'}f_U,\qquad g_T=g_{T'}f_T
\]
% archive sheet number: 118
\scanpage{119}
% line head | text | 68
\mspage{68}
% line 1–2 | text | Donc la donnée d'un triple
Donc la \deleted{\unsure{coup.}}\inserted{donnée d'un} triple $(*)$ équivaut à celle d'un
triple
% line 3–5 | display | (g_U prime,g_T prime,g_X)
\[
\begin{aligned}
(g_{U'},g_{T'},g_X)\ \big|\ &g_X\text{ une flèche }X\to X''\text{ correspondant}\\
&\text{à }g_{\unsure{U|T}}=g_{U'}\circ f_U\text{ et }g_T=g_{T'}\circ f_T\text{ déjà}\\
&\deleted{\text{\unsure{déterminés}}}\ \text{\inserted{\unsure{donné|donnés}} (via }g_{U'},g_{T'}\text{),}
\end{aligned}
\]
% line 6 | text | donc en vertu de 1.3.
donc en vertu de 1.3. cela correspond à
% line 7–10 | display | (*) rewritten triple, with evaluation under the bracket
\[
\begin{aligned}
(g_{U'},g_{T'},u)\quad\left|\begin{aligned}
&g_{U'}:U'\to X''\\
&g_{T'}:T'\unsure{-|\to}X''\\
&u:H\longrightarrow
\underset{\scriptstyle (u,t)\longmapsto\operatorname{Hom}_{X''}(g_{U'}f_U(u),g_{T'}f_T(t))}
{\underline{H_{X''}\bigl((g_{U'}\circ f_U)^{\mathrm o}\times(g_{T'}\circ f_T)}} .
\end{aligned}\right.
\end{aligned}\tag*{(*)}
\]
% line 11–12 | text | On va interpréter la donnée de u
On va interpréter la donnée de $u$, en termes d'un
foncteur
% line 13 | display | (**) F upper star
\[
F^*=(H'\mapsto H'\circ F):
\underline{\operatorname{Hom}}({U'}^{\mathrm o}\times T',\mathrm{Ens})\to
\underline{\operatorname{Hom}}(U^{\mathrm o}\times T,\mathrm{Ens})\tag*{(**)}
\]
% line 14 | text | où F is given by
où $F:U^{\mathrm o}\times T\longrightarrow {U'}^{\mathrm o}\times T'$ est donné par $F=f_U^{\mathrm o}\times f_T$.

% line 15–16 | text | On a un objet H double prime
On a un objet $H''$ \inserted{$=H''_{g_{U'},g_{T'}}$} de la catégorie source de $(**)$
déterminé \deleted{par} en termes du couple $(g_{U'},g_{T'})$, par
% line 17 | display | H double prime evaluated at u prime and t prime
\[
H''(u',t')=\operatorname{Hom}_{X''}(g_{U'}(u'),g_{T'}(t')),
\]
% line 18–19 | text | et le bifoncteur but
et le bifoncteur but de la flèche $u$ dans $(*)$ s'interprète comme
% line 20–21 | display | F upper star of H double prime
\[
\begin{aligned}
F^*(H'')=H''\circ F:(u,t)\unsure{\vdash|\mapsto}{}&
 H''(f_U(u),f_T(t))\deleted{\unsure{=}}\\
&\underset{\operatorname{Hom}(g_{U'}f_U(u),g_{T'}f_T(\unsure{t|t'}))}{\shortparallel}.
\end{aligned}
\]
% line 22 | text | Ainsi, u s'interprète
Ainsi, $u$ s'interprète comme une flèche
% line 23 | display | u to F upper star of H double prime
\[
u:H\longrightarrow F^*(H'').
\]
% line 24–25 | text | Or le foncteur F upper star
Or le foncteur $F^*$ a un adjoint à gauche
$F_!$, \deleted{\unsure{et}}\inserted{i.e.} on a
% line 26 | display | (***) Hom adjunction
\[
\operatorname{Hom}(H,F^*(H''))\simeq\operatorname{Hom}(F_!(H),H'')\tag*{(***)}
\]
% archive sheet number: 119
\scanpage{120}
% line head | text | 69
\mspage{69}
% line 1 | text | avec
avec
% line 2 | display | F lower shriek between the two Hom categories
\[
F_!:\underline{\operatorname{Hom}}(U^{\mathrm o}\times T,\mathrm{Ens})\longrightarrow
\underline{\operatorname{Hom}}({U'}^{\mathrm o}\times T',\mathrm{Ens}).
\]

% line 3 | text | Posons
\deleted{L'objet} Posons
% line 4 | display | H prime equals F shriek of H
\[
H'=F_!(H):{U'}^{\mathrm o}\times T'\to\mathrm{Ens},
\]
% line 5–7 | text | c'est un bifoncteur bien déterminé
c'est un bifoncteur bien déterminé \deleted{\unsure{étant}} \unsure{la}
seule donnée du diagramme initial \inserted{$(\unsure{*|\alpha})$ p. 67}, ou
ce qui revient au $=$, en termes de
% line 8 | display | U,T,H,f_U,f_T
\[
(U,T,H,f_U:U\unsure{-|\to}U',f_T:T\to T').
\]
% line 9–11 | text | Et on voit que
Et on voit que la donnée d'une flèche
$X^{\unsure{\prime|\prime\prime}}\longrightarrow X''$, $X''$ objet variable de $\mathrm{Cat}$, équivaut
à celle de
% line 12 | display | the two g's and u
\[
g_{U'}:U'-X'',\quad g_{T'}:T'\to X'',\quad
u:H'\to H''=H_{X''}(g_{U'}^{\mathrm o}\times \unsure{g_{T'}|g_{T'}^{\mathrm o}|f_{T'}}),
\]
% line 13–14 | text | c'est à dire, en vertu de 1.3
c'est à dire, en vertu de 1.3, à celle d'une
flèche
% line 14 right | display | X prime subscript nought
\[
X'_0\longrightarrow X'',
\]
% line 15–19 | text | où X prime nought désigne
où $X'_0$ désigne l'objet de $\mathrm{Cat}$ déduit
de \unsure{(}$U',T'$ par recollement au moyen de $H'$.
Cela montre donc que la limite \unsure{inductive}
cherchée n'est autre que $X'_0$. On a
ainsi trouvé le

% line 20–23 | text | Théorème 1. La limite inductive
\uline{Théorème 1}\quad La \deleted{\unsure{somme amalgamée}}\inserted{limite inductive} \inserted{$X'$} dans $\mathrm{Cat}$ du
diagramme $(\alpha)$ page 67, défini en termes
d'un triple $(X,U,T)$ de $\unsure{\mathscr T|\mathcal T}$ et de flèches
$f_U:U-U'$, $f_T:T\to T'$ dans $\mathrm{Cat}$, existe (\unsure{sans mérite|cas mérite|cas mixte})
% archive sheet number: 120
\scanpage{121}
% line head | text | 70
\mspage{70}
% line 1 | text | et les flèches canoniques
et les flèches canoniques
% line 2 | diagram | U prime, X prime, T prime
\diagram{\begin{tikzcd}[ampersand replacement=\&]
U'\arrow[r,hook,"i'"] \& X' \& T'\arrow[l,hook',"j'"']
\end{tikzcd}}
% line 3–5 | text | sont des immers. ouverte et fermée
sont des immers. \unsure{ouverte|ouvertes} et \unsure{fermées|fermée} complémentaires
l'une de l'autre. La donnée de recollement
qui \unsure{définit} $X'$ via $U',T'$ est donnée par
% line 6 | display | H prime in terms of the functor F
\[
H'=(\underbrace{f_U^{\mathrm o}\times f_T}_{F})_!(H),:
{U'}^{\mathrm o}\times {T'}^{\deleted{\unsure{\mathrm o}}}\unsure{-|\to}\mathrm{Ens},
\]
% line 7–9 | text | où H est la donnée de recollement
où $H$ est la donnée de recollement qui définit
$X$, i.e. $H\deleted{\unsure{(u,t)}}:U^{\mathrm o}\times T-\mathrm{Ens}$ est le bifoncteur
$(u,t)\mapsto\operatorname{Hom}_X(u,t)$.

% line 10 | text | Ainsi on a, pour
Ainsi on a, pour $u'$ dans $U'$ et $t'$ dans $T'$
% line 11 | display | Hom in X prime
\[
\operatorname{Hom}_{X'}(u',t')=\deleted{\unsure{\lim_{\rightarrow}}_{U^{\mathrm o}\times T}}\ F_!(H)(u',t'),
\]
% line 12–15 | text | et il reste à expliciter
et il reste à expliciter $F_!(H)$. Ici \deleted{$F$ est un}
$F=f_U^{\mathrm o}\times f_T$ est un foncteur \unsure{covariant} sur $U^{\mathrm o}\times T$,
on peut l'interpréter comme un préfaisceau sur
\deleted{$\unsure{U}$} $U\times T^{\mathrm o}$, et écrire dès lors
% line 16–17 | display | (*) Kan-extension value, indexing category below lim
\[
\underbrace{F_!(H)(u',t')}_{\operatorname{Hom}_{X'}(u',t')}
=\underset{{}_{(u',t')}\backslash^{U\times T^{\mathrm o}}}{\lim_{\rightarrow}}
\underbrace{\operatorname{Hom}_X(u,t)}_{H(u,t)}.\tag*{(*)}
\]
% line 18–20 | text | La catégorie sur laquelle
La catégorie \inserted{$\mathord{}_{(u',t')}\backslash^{U\times T^{\mathrm o}}$} sur laquelle on prend la $\lim_{\rightarrow}$
\deleted{est donc celle des couples de flèches}
s'explicite comme un produit
% line 21 | display | factorization of the indexing category
\[
{}_{(u',t')}\backslash^{U\times T^{\mathrm o}}\simeq
({}_{u'}\backslash^{U})\times(T_{/t'})^{\mathrm o}.
\]
% line 22 | text | Les applications
Les applications
% line 23 | display | Hom_X to Hom_X prime
\[
\operatorname{Hom}_X(u,t)\longrightarrow\operatorname{Hom}_{X'}(u',t')
\quad(=F_!\unsure{(H]|(H)}(u',t'\unsure{]|)}),
\]
% archive sheet number: 121
\scanpage{122}
% line head | text | 71
\mspage{71}
% line 1–3 | text | déduites de cette formule
déduites de cette formule ne sont autres que \deleted{celles}
les suivantes : $u,t$ sont donnés comme objets de
$\mathord{}_{u'}\backslash^{U}$, $T_{/t'}$, i.e. avec des flèches
% line 4 | display | alpha and beta
\[
u'\xrightarrow{\alpha}f_U(u),\qquad
\deleted{\unsure{f}}f_T(t)\mathrel{\overset{\beta}{\unsure{-|\to}}}t'
\]
% line 5 | text | d'où
d'où
% line 6 | diagram | Hom_X via f, then via alpha and beta
\diagram{\begin{tikzcd}[ampersand replacement=\&,column sep=small]
\operatorname{Hom}_X(u,t)
\arrow[r,"{\text{via }\deleted{\unsure{g}}f=f_X:X\to X'}"]
\& \operatorname{Hom}_{X'}(\underset{f_U(u)}{\underline{f(u)}},\underset{f_T(t)}{\underline{f(t)}})
\arrow[r,"{\text{via }\alpha{,}\beta}"]
\& \operatorname{Hom}_{X'}(u',t').
\end{tikzcd}}

% line 7–9 | text | Notons pour mémoire
Notons pour mémoire que pour tout
objet $Z$ de $\mathrm{Cat}$, et tout préfaisceau $G$
sur $Z$, on a la formule
% line 10 | display | lim G and pi nought
\[
\underset{Z}{\lim_{\rightarrow}}G\simeq\pi_0(Z_{/G}),
\]
% line 11–23 | text | où Z/G s'interprète
où $Z_{/G}$ s'interprète en terme\unsure{s} de l'immersion
canonique $Z\longrightarrow Z^\wedge$ ($\unsure{C|G}\in\operatorname{Ob}Z^\wedge$).
C'est donc une catégorie \deleted{\unsure{co}}fibrée sur $Z$ à fibres
discrètes, \unsure{\{}et on a \inserted{en fait} un \deleted{équivalence}\inserted{isom.} entre la
\inserted{ss} catégorie \inserted{pleine} de $\unsure{\mathrm{Kat}|\mathrm{Cat}}_{/Z}$ formée \unsure{de|des} telles catégories
fibrées à fibres discrètes, avec $Z^\wedge\unsure{\}|\big)}$.
Dans ce dictionnaire parfait entre préfaisceaux
d'une\unsure{.}, et cat. fibrées à fibres discrètes,
le foncteur $\lim_{\rightarrow}$ sur $Z^\wedge$ correspond simplement
au \inserted{foncteur} $\pi_0$ \deleted{\unsure{dans}} de la catégorie des objets au
dessus de $Z$ dans $\mathrm{Cat}$, lui-même induit
par le foncteur (indépendant de $Z$)
% line 24 | display | Cat to (Ens), pi nought
\[
\mathrm{Cat}\xrightarrow{\pi_0}(\mathrm{Ens}),
\]
% archive sheet number: 122
\scanpage{123}
% line head | text | 72
\mspage{72}
% line 1–3 | text | On peut donc réécrire
On peut donc réécrire la formule fondamentale
$(*)$ p. 70, en termes du foncteur $H$ définissant
le recollement $X$ de $U$ et de $T$, par
% line 4–5 | display | (*) boxed component formula and two identifications
\[
\boxed{\underbrace{\operatorname{Hom}_{X'}(u',t')}_{H'(u',t')}
\simeq\pi_0\bigl(\underbrace{{}_{u',t'}\backslash^{U\times T^{\mathrm o}}_{/H}}
_{({}_{u'}\backslash^{U}\times(T_{/t'})^{\mathrm o})_{/H}}\bigr)}\tag*{(*)}
\]

% line 6–11 | text | Pour la suite
\unsure{Pour|Par} la suite, nous n'aurons sans doute qu'à
utiliser le cas où on a soit $T'=T$, $f_T=\mathrm{id}_T$,
soit l'inverse. Nous travaillerons surtout dans
le premier cas, en déduisant ce dont nous
aurons besoin pour le cas dual, par principe de
dualité (\unsure{interne} $X\unsure{\vdash|\mapsto}X^{\mathrm o}$ de la catégorie $\mathrm{Cat}$).
% line 12–14 | text | Dans ce cas, la formule fondamentale
Dans ce cas, \inserted{\unsure{ainsi}} la formule fondamentale $(*)$, \deleted{\unsure{de}} l'élmt
$t'$ de $T'$ \unsure{s'écrit|s'écrivent} comme un élément $t$ de $T$,
et la formule se retranscrit comme
% line 15–17 | display | (**) boxed special case and category of triples
\[
\boxed{\operatorname{Hom}_{X'}(u',t)\simeq
\pi_0\bigl(\underset{\substack{\text{catégorie des}\\\text{triples }(u,\alpha,\beta'),}}{\underline{{}_{u'}\backslash^{U}_{/t}}}\bigr)}
\ \left\{\begin{aligned}
&u\in\operatorname{Ob}U\\
&\alpha:u\to t\\
&\beta':u'\to f_U(u).
\end{aligned}\right.\tag*{(**)}
\]

% line 18–23 | text | C'est la formule que j'avais trouvée
C'est la formule que j'avais trouvée directement
dans ce cas, sans me soucier \inserted{encore} de faire varier $T$
en même temps que $U$. Je voudrais faire le
lien entre les deux formules $(*)$ et $(**)$, par
voie topologique, \unsure{donc} établir une bijection
canonique, fonctorielle etc
% line 24 | display | desired bijection between two pi noughts
\[
\pi_0({}_{u'}\backslash^{U}_{/t})\mathrel{\overset{\sim}{\longrightarrow}}
\pi_0\bigl(({}_{u'}\backslash^{U}\times(T_{/t'})^{\mathrm o})_{/H}\bigr).
\]
% archive sheet number: 123
\scanpage{124}
% line head | text | 73
\mspage{73}
% line 1 | text | On a sûrement un foncteur
On a sûrement un foncteur
% line 2 | display | Comparison functor
\[
{}_{u'}\backslash^{U}_{/t}\longrightarrow
\bigl({}_{u'}\backslash^{\unsure{U'|U}}\times(T_{/t})^{\mathrm o}\bigr)_{/H}
\]
% line 3 | text | qui associe à un triple
qui associe à un triple
% line 4 | display | (u, alpha, beta prime)
\[
(u,\alpha,\beta'):\quad u\in\operatorname{Ob}U,\quad\alpha:u\unsure{-|\to}t,\quad
\beta':u'\to f_U(u)
\]
% line 5–7 | text | le triple
le triple $(\bar u,\bar t,\unsure{\xi|\zeta})$ ($\xi\in H(\bar u,\bar t)$) \deleted{du \unsure{sec.}} dans
le second membre, avec $\bar u$ dans $\mathord{}_{u'}\backslash^{U'}$, $\bar t$ dans $T_{/t}$,
défini ainsi :
% line 8 | display | barred u, barred t, xi
\[
\bar u=(u,\beta'),\quad \bar t=(t,\mathrm{id}_t),\quad
\underset{\substack{\shortparallel\\\alpha}}{\xi}\in H(\bar u,\bar t)=\operatorname{Hom}(u,t).
\]

% line 9 | text | En somme
En somme,
% line 10 | display | Fibre at the initial object
\[
\begin{aligned}
{}_{u'}\backslash^{U}_{/t}&=\text{fibre de }\bigl({}_{u'}\backslash^{U}\times(T_{/t})^{\mathrm o}\bigr)_{/\unsure{H|t}}
\text{ en l'objet \deleted{final}\inserted{initial} }\bar t\\
&\hspace{17em}\text{de }(T_{/t})^{\unsure{\mathrm o|'}}
\end{aligned}
\]
% line 11 | text | par la projection naturelle
par la projection naturelle, flèche en
% line 12–14 | diagram | three levels of the natural projection
\diagram{\begin{tikzcd}[ampersand replacement=\&,row sep=small]
\bigl({}_{u'}\backslash^{U}\times(T_{/t})^{\mathrm o}\bigr)_{/H}\arrow[d]\\
{}_{u'}\backslash^{U}\times(T_{/t})^{\mathrm o}\arrow[d,"\pi"]\\
(T_{/t})^{\mathrm o}
\end{tikzcd}}
% line 15–21 | text | Il faut voir que l'inclusion
Il faut voir que l'inclusion de la fibre donne
une bijection sur les $\pi_0$, ou ce qui revient
au $=$, que l'inclusion de la fibre
$\mathord{}_{u'}\backslash^{U}$ \uline{de $\pi$} en l'objet initial $(t,\mathrm{id}_t)$ de $(T_{/t})^{\mathrm o}$,
donne une bijection sur les $\lim_{\rightarrow}$ \unsure{du|des}
\unsure{foncteur|foncteurs} $H$ !\unsure{!} Donc l'inclusion a intérêt
à être asphérique, ou du moins co-cofinal.
% archive sheet number: 124
\scanpage{125}
% line head | text | 74
\mspage{74}
% line 1–5 | text | Mais l'inclusion d'un objet initial
Mais l'inclusion d'un objet initial est toujours
asphérique, et le reste par produit avec un
objet quelconque (ici $\mathord{}_{u'}\backslash^{U}$) de $\mathrm{Cat}$
(et $=$ par \unsure{tt} changement de base lisse\ldots),
d'où la conclusion. \unsure{$\backslash$}

% line 6–9 | text | Dans la suite, nous nous intéressons
Dans la suite, nous nous intéressons
\unsure{surtout|seulement} à la situation de somme amalgamée
(où on change de notations un peu, à
présent)
% line 10–12 | diagram | (*) two rows of the pushout with dotted closed sides
\diagram{\begin{tikzcd}[ampersand replacement=\&,column sep=small]
 \& U\arrow[r,hook,"i"]\arrow[d,"p"'] \& X\arrow[d,"q"'] \& T\arrow[l,dashed,hook',"j"']\arrow[d,dashed]\\
 {(*)} \& U'\arrow[r,hook,"i'"] \& X' \& {T'\simeq T}\arrow[l,dashed,hook']
\end{tikzcd}}
% line 13–22 | text | Le fermé T prime complémentaire
Le \deleted{catégorie}\inserted{fermé $T'$} complémentaire de $U'$ \inserted{dans $X'$} est alors
can. isom. au fermé $T$ compl. de $U$ dans
$X$, l'isom. étant induit par $q$. Souvent
nous identifierons $T'$ et $T$, \inserted{ce qui est licite,} en regardant
$X'$ comme déduit du recollement de
\deleted{$\operatorname{Ob} X'=$}
$U'$ et de $\unsure{t|T}$, par le bifoncteur $\mathord{U'}^{\mathrm o}\times\unsure{\bar T|T}\unsure{-|\to}\mathrm{Ens}$
explicité dans la formule $(**)$ p. 72. Nous
nous \unsure{intéressons|intéresserons} surtout à dégager des conditions
qui assurent les deux implications suivantes
% line 23–24 | display | a) and b), membership in W
\[
\begin{aligned}
\text{a)}\quad i\in W&\Longrightarrow\unsure{,}\ i'\in W\quad(\text{et }=\ \Longleftrightarrow)\\
\text{b)}\quad p\in W&\Longrightarrow q\in W\quad(\text{et }=\ \Longleftrightarrow).
\end{aligned}
\]
% archive sheet number: 125
\scanpage{126}
% line head | text | 75
\mspage{75}
% line 1–3 | text | Prop 4
\uline{Prop 4} : Soit $t\in\inserted{\operatorname{Ob}}T$. Pour que $t\in\inserted{\operatorname{Ob}}\bar U$, il f. et s. que
$q(t)\in\inserted{\operatorname{Ob}}\overline{U'}$. En d'autres termes, le ``bord'' $\dot T=T\cap\bar U$ de $T$
dans $X$ est ``le même'' que son bord dans $X'$.

% margin note written sideways
\margin{\uline{Proposition 5} : Soit
\diagram{\begin{tikzcd}[ampersand replacement=\&,row sep=tiny,column sep=tiny]
Y\arrow[r,"i"]\arrow[d,"p"'] \& X\arrow[d,"q"]\\
Y'\arrow[r,"i'"'] \& X'
\end{tikzcd}}
un \deleted{\unsure{ca}}carré cocartésien dans $\mathrm{Cat}$. Si $i$ est une immersion ouverte (resp. fermée) alors $i'$ aussi, \inserted{et le carré est aussi cartésien.} \unsure{idem|id.} pour $p$ et $q$. Si $i$ et $p$ sont des imm. ouvertes \inserted{(resp. \unsure{fermées})} \deleted{\unsure{alors}} la situation est isom. à celle d'un objet $X'$ de $\mathrm{Cat}$, de deux ouverts (resp. fermés) $X,Y'$ de $X$ ay. \unsure{pr|pour} \unsure{réunion|somme} $X'$, pour intersection $Y$. Inversement, dans une telle situation, \deleted{$X$ s'id} le carré correspondant est \unsure{cocartésien|cartésien}, i.e. $X$ s'identifie à la somme amalgamée
$X'\simeq X\amalg_Y Y'$.
\unsure{Nous retrouvons|Pour ceci cf} \unsure{en effet|au §} la prop 15, page 158, et son cor. \unsure{Ils disent}}
% line 4 | text | Dém. Dire que
Dém. Dire que $t\in\operatorname{Ob}\bar U$ signifie aussi
% line 5 | display | U over t nonempty
\[
U_{/t}\ne\varnothing
\]
% line 6–8 | text | et implique, par continuité
et implique, par continuité, que $q(t)\in\deleted{\unsure{\operatorname{Ob}}}\inserted{\overline{q(U)}\subset}\overline{U'}$. Inversement,
si $q(t)\in\overline{U'}$, i.e. s'il existe $u'\in U'$
tel que
% line 9 | display | Hom(u prime,q(t)) nonempty
\[
\deleted{\operatorname{Hom}(u'}\quad\operatorname{Hom}(u',q(t))
\bigl(=\pi_0({}_{u'}\backslash^{U}_{/t})\bigr)\ne\varnothing,
\]
% line 10–11 | text | en résulte que
en résulte que $\mathord{}_{u'}\backslash^{U}_{/t}\ne\varnothing$, a fortiori
$U_{/t}$ (qui \unsure{reçoit} $\mathord{}_{u'}\backslash^{U}_{/t}$) $\ne\varnothing$, \unsure{qed}
\margin{\uline{NB} Nous utilisons ``\uline{immersion ouverte}'', ``\uline{immersion fermée}'' au \uline{sens strict} : c'est un \uline{monomorphisme} \unsure{pleinement|entièrement} fidèle dont l'image \deleted{\unsure{essentielle}} est une sous-catégorie \uline{pleine} \unsure{strictement} ouverte ou fermée.}

% line 12 | text | 2. Sommes amalgamées et W-équivalences
2\unsure{.}\quad \uline{Sommes amalgamées et $W$-équivalences}.

% line 13 | text | 2.1. Axiome du carré cocartésien
2.1. \uline{Axiome du carré cocartésien et axiome du cube} (ou ``\uline{lemme des cinq}'')

% line 14–15 | text | Nous supposons à présent donné
Nous supposons à présent donné un localiseur
fondamental
% line 16 | display | W contained in Fl(Cat)
\[
W\subset\operatorname{Fl}(\mathrm{Cat}),
\]
% line 17–24 | text | qui satisfait à W(1,2bis,3)
qui satisfait à $W(1,2\mathrm{bis},3)$. Nous ne supposons
pas pour le moment \unsure{pour|que} les \deleted{\unsure{autres}} autres axiomes
\unsure{suivants|qu'il ait} $W(4.)$ \deleted{$\unsure{W(5)}$} (axiome de fibration), $W(5)$
(axiome de \deleted{limites inductives}\inserted{\unsure{connexité}}), $W(6)$ (axiome des
limites inductives). Mais nous allons introduire
à présent deux axiomes \inserted{de stabilité} supplémentaires, de
nature différente, \deleted{\unsure{car ils}} concernant les sommes
amalgamées.

% line 25–28 | text | W(7bis), Axiome du carré cocartésien
\uuline{$W(7\mathrm{bis})$} \inserted{(Axiome du carré cocartésien)} Soient $X$ un objet de $\mathrm{Cat}$, $U,U'$ deux
ouverts de $X$ tel que $U\cup U'=X$, considérons le
diagramme cocartésien correspondant (cf prop. 5 en
\deleted{\unsure{marge}} de la page) :
% archive sheet number: 126
\scanpage{127}
% line head | text | 76
\mspage{76}
% line 1–3 | diagram | (*) intersection, U, U prime, X
\diagram{\begin{tikzcd}[ampersand replacement=\&,row sep=small,column sep=small]
\& \& U\cap U'\arrow[dl,hook',"i'_0"']\arrow[dr,hook,"i_0"] \&\\
{(\unsure{*|\alpha})} \& U\arrow[dr,hook,"i"'] \& \& U'\arrow[dl,hook',"i'"]\\
\& \& X \&
\end{tikzcd}}
% line 4 | text | Alors
Alors
% line 5 | display | i_0 in W
\[
i_0\in W\Longleftrightarrow i\in W
\]
% line 6 | text | et symétriquement
(et symétriquement
% line 7 | display | i_0 prime in W
\[
i'_0\in W\Longleftrightarrow i'\in W\ \bigr).
\]
% line 8–9 | text | NB Par dualité
\uline{NB} Par dualité cet axiome est équivalent à l'axiome similaire
avec des \uline{fermés} $T,T'$ de $X$ de réunion $X$.

% line 10–11 | text | Corollaire
\uline{Corollaire}\quad \inserted{\fbox{Supposons $W(7\mathrm{bis})$.}} Considérons les conditions \unsure{suiv.} \deleted{\unsure{de}} sur
le carré précédent

% line 12 | text | a) U intersection U prime
a) $U\cap U'$ \inserted{$W$-}asphérique

% line 13 | text | b) U, U prime
b) $U,U'$ $W$-asphériques

% line 14 | text | c) X
c) \unsure{$X$} $W$-asphérique

% line 15 | text | Alors (a et b)
Alors (a et b) \unsure{impliquent} c, et (b et c) implique a.

% line 16–22 | text | NB Moralement
\unsure{\{} \uline{NB} Moralement, il faudrait \inserted{aussi} que (c \deleted{$+$} et a) implique
\deleted{\unsure{aussi}} b, i.e. si $U\cap U'$ et $X$ sont $W$-asphériques,
il en est de même de $U$ et de $U'$. Mais
il ne semble pas que l'axiome du carré cocartésien
permette de le prouver. Par contre, l'axiome
qui suit impliquera le corollaire précédent sous sa
forme complète.

% line 23–24 | text | W, Axiome du cube, Lemme des cinq
\deleted{\uline{\unsure{NB}}} $W(\unsure{7|8})$ (\inserted{``Axiome du cube'', ou} ``Lemme des cinq'') Soit un diagramme
commutatif cubique
% line 25–28 left | diagram | cancelled drawing, intersection at top through X1, lower part begun
\diagram{\deleted{\begin{tikzcd}[ampersand replacement=\&,row sep=tiny,column sep=tiny,execute at end picture={
\draw[thin] (current bounding box.north) -- (current bounding box.west) -- (current bounding box.north east) -- (current bounding box.south west) -- (current bounding box.east) -- (current bounding box.south);
}]
\& \& U_1\cap U'_1\arrow[dl,hook',"i'_0"']\arrow[dr,hook,"i_0"] \&\\
\& U_1\arrow[dr,"i"']\arrow[ddl] \& \& U'_1\arrow[dl,"i'"]\\
\& U_2\cap U'_2\arrow[dl] \& X_1 \&\\
U_2 \& \& \&
\end{tikzcd}}}
% line 25–28 right | diagram | live cube comparing intersections, opens and unions
\diagram{\begin{tikzcd}[ampersand replacement=\&,row sep=small,column sep=small]
U_1\cap U'_1\arrow[rr]\arrow[dr]\arrow[dd,"f_{U_1\cap U'_1}"'] \& \& U_1\arrow[dr]\arrow[dd,"f_{U_1}" description] \&\\
\& U'_1\arrow[rr]\arrow[dd,"f_{U'_1}"] \& \& X_1\arrow[dd,"f"]\\
U_2\cap U'_2\arrow[rr]\arrow[dr] \& \& U_2\arrow[dr] \&\\
\& U'_2\arrow[rr] \& \& X_2
\end{tikzcd}}
% line 29 | text | représentant un morphisme
représentant un morphisme (flèches verticales) d'un
% archive sheet number: 127
\scanpage{128}
% line head | text | 77
\mspage{77}
% line 1 | text | carré cocartésien
carré cocartésien
% line 1–2 | diagram | intersection, U1, U1 prime, X1
\diagram{\begin{tikzcd}[ampersand replacement=\&,row sep=small,column sep=small]
U_1\cap U'_1\arrow[r]\arrow[d] \& U_1\arrow[d]\\
U'_1\arrow[r] \& X_1
\end{tikzcd}}
% line 1–4 | text | du type envisagé
du type envisagé dans $W(7)$ ($X_1$ réunion de deux ouverts $U_1,U'_1$) dans un
autre (\inserted{d'}indices $2$). Considérons les conditions suivantes

% line 5 | text | a) f on the intersection
a) $f_{U_1\cap U'_1}\in W$

% line 6 | text | b) f on the opens
b) $f_{U_1}$ et $f_{U'_1}$ sont $\in W$

% line 7 | text | c) f
c) $f\in W$.

% line 8–11 | text | Alors deux parmi ces conditions
Alors deux parmi ces conditions impliquent la troisième.
En d'autres termes, quand l'une des \unsure{ces|les} conditions est satisfaite, les deux autres sont équivalentes
entre elles.

% line 12–14 | text | NB Cet axiome
\uline{NB} Cet axiome est équivalent à l'axiome dual,
\deleted{\unsure{dans}} en termes d'objets $X_1,X_2$ représentés \inserted{chacun} comme
réunion de deux fermés.

% line 15–18 | text | Corollaire
\uline{Corollaire} \deleted{\unsure{Comme|Comment}} \inserted{Supposons $W(8)$ satisfait. Alors} la conclusion \unsure{de|à} $W(7\mathrm{bis})$ (axiome
qu'on n'a pas à utiliser ici), sous sa forme complétée
(cf le NB qui le suit) est valable: deux parmi les
conditions a), b), c) impliquent la troisième.

% line 19–21 | text | Il suffit de regarder
Il suffit de regarder l'unique morphisme
du diagramme $(\unsure{\alpha|a})$ p.~76, dans le diagramme
trivial
% line 21 | diagram | four terminal categories e
\diagram{\begin{tikzcd}[ampersand replacement=\&,row sep=tiny,column sep=tiny]
\& e\arrow[dl]\arrow[dr] \&\\
e\arrow[dr] \& \& e\arrow[dl]\\
\& e \&
\end{tikzcd}}
% line 21–24 | text | correspondant au cas où
, correspondant au cas où
$X$ est la catégorie finale $e$, \inserted{et} $U',U''$ égaux à $X$. Alors le
corollaire n'est autre que $W(8)$ appliqué à ce
cas particulier.

% line 25–27 | text | Cet argument fait sentir
Cet argument fait sentir que $W(\unsure{7|8})$ est \inserted{sans doute} plus
fort que $W(7\mathrm{bis})$. Incité par ce sentiment, je \unsure{me|m'}
aperçois avec plaisir qu'il est fondé:
% archive sheet number: 128
\scanpage{129}
% line head | text | 78
\mspage{78}
% line 1–2 | text | Proposition 6
\uline{Proposition 6} L'axiome $W(\unsure{7|8})$ (lemme des cinq)
implique $W(7\mathrm{bis})$ (axiome du carré cocartésien).

% line 3 | text | On regarde le diagramme
On regarde le diagramme cubique suivant
% line 4–9 | diagram | cube from intersection, opens and X to repeated U prime and X
\diagram{\begin{tikzcd}[ampersand replacement=\&,row sep=large,column sep=small]
U\cap U'\arrow[rr,"i_0^{\deleted{\unsure{\prime}}}"]\arrow[dr,"i'_0","\deleted{\unsure{i_0}}"']\arrow[dd,"f_{U\cap U'}=i_0"'] \& \& U'\arrow[dr,"i'"]\arrow[dd,"f_{U'}=\mathrm{id}" description] \&\\
\& U\arrow[rr,"i"]\arrow[dd,"f_U=i"'] \& \& X\arrow[dd,"f=\mathrm{id}_X"]\\
U'\arrow[rr,"\mathrm{id}"]\arrow[dr,"i'"'] \& \& U'\arrow[dr,"i'"] \&\\
\& X\arrow[rr,"\mathrm{id}"] \& \& X
\end{tikzcd}}
% line 10–16 | text | du carré cocartésien donné
, du carré cocartésien donné dans le carré cocartésien
évident de $X$ et de ses deux ouverts $X,U'$,
d'intersection $U'$. Ici les deux flèches verticales
$f$ et $f_{U'}$ sont des iso donc \deleted{sont} $\in W$. Avec les
notations de $W(7)$, on a donc c) ($f\in W$ \unsure{!}),
donc a) et b) sont équivalentes, i.e. $i_0\in W$
et $\unsure{\{|(}i\in W$ sont équivalentes \hfill \unsure{qed.}

% line 17–18 | text | Théorème 2
\uline{Théorème 2} Le localiseur fond. $W_{\infty}$ satisfait à
l'axiome $W(7)$.

% line 19–24 | text | C'est sûrement une conséquence
C'est \inserted{sûrement} une conséquence \inserted{du critère \unsure{cohomologique|homotopique} d'Artin-Mazur, compte tenu} \unsure{des|de} la suite exacte
de \unsure{Mayer|Meyer}-Vietoris pour l'un \deleted{\unsure{deux|des}} et l'autre
des deux carrés cocartésiens, et les variantes
non commutatives de \unsure{Meyer|Mayer}-Vietoris. Je n'ai
pas fait la vérification en forme -- et devrais
sans doute réparer cette omission.
% archive sheet number: 129
